\documentclass[a4paper,12pt]{ctexart}
\usepackage{xcolor,graphicx,geometry,array,hyperref,booktabs}
\hypersetup{colorlinks=true,linkcolor=blue!70!black}
\geometry{top=1.8cm,bottom=1.8cm,left=1.8cm,right=1.8cm}
\linespread{1.35}
\newenvironment{thinkbox}{\vspace{0.3cm}\begin{center}\begin{minipage}{0.88\textwidth}\colorbox{blue!5}{\begin{minipage}{0.94\textwidth}\vspace{0.15cm}\textbf{【想一想】}}\vspace{0.15cm}\end{minipage}\end{minipage}\end{center}\vspace{0.3cm}}{}
\newenvironment{funfact}{\vspace{0.2cm}\begin{center}\begin{minipage}{0.88\textwidth}\fbox{\begin{minipage}{0.92\textwidth}\vspace{0.1cm}\textbf{【有趣的小知识】}}\vspace{0.1cm}\end{minipage}\end{minipage}\end{center}\vspace{0.2cm}}{}
\newenvironment{figplaceholder}[1]{\vspace{0.2cm}\begin{center}\fbox{\begin{minipage}{0.82\textwidth}\centering\textbf{【插图位置】}\\[0.2cm]#1\end{minipage}}\end{center}\vspace{0.2cm}}{}

\title{\Huge\textbf{土木工程小故事}}
\date{}
\begin{document}\maketitle\thispagestyle{empty}

\section{桥为什么不会塌？——从赵州桥到港珠澳大桥}

小明和爸爸开车经过一座高架桥——桥下的路离他们越来越远，"爸，这桥这么高，为什么不会塌？"

爸爸说："你看桥下面的结构——很多三角形组成的钢架，三角形是不可变形的。一座桥的受力路径是这样的：桥面上的荷载→传到主梁→再传到墩柱→最后传到地基。每一步都把力分散和传递到了更大的范围。"

世界上有几种主要的桥型：\textbf{梁桥}、\textbf{拱桥}、\textbf{斜拉桥}和\textbf{悬索桥}。拱桥把重力转化为向两侧推的压力（最适合用石材——石头抗压很强）；悬索桥把力传到主缆上（适合超大跨度）。

\begin{figplaceholder}
插图：四种桥型的受力示意——梁桥（弯矩）、拱桥（压力）、斜拉桥（斜索拉力）、悬索桥（主缆拉力）
\end{figplaceholder}

\textbf{赵州桥：}建于公元605年（隋朝），跨度37米，由工匠李春设计——1400多年来依然坚挺。欧洲类似的拱桥直到14世纪才出现（晚了700年）。赵州桥敞肩式拱形设计用石头之间的相互挤压来传力——没用任何黏合剂。

\textbf{港珠澳大桥：}全长55公里（世界最长跨海大桥），设计使用寿命120年。需要应对的问题：台风（风速最高可达16级）、海水腐蚀（混凝土和钢材的耐久性）、地震（位于华南地震带）、船舶撞击（主航道桥可承受30万吨船只的撞击）。

\begin{thinkbox}
你家附近有哪些桥？去仔细观察它们的结构——是梁桥、拱桥还是斜拉桥？桥墩是怎么连到河床的？
\end{thinkbox}

\begin{funfact}
北京的三元桥在2015年进行了一次"空中换桥"——整座桥被上千个千斤顶同步顶升，然后平移替换，施工过程只中断了43小时交通。这是国内首例大型桥梁的"整体置换"。
\end{funfact}

\newpage

\section{北京最高的楼怎么建起来的？——地基与抗震}

国贸三期高330米——如果把它放倒在路边，它几乎有长安街那么长。这么高的楼，没有被风吹倒、没有被地震震塌——秘密在两方面：

\textbf{地基：}摩天大楼的地基不是"放在地面上"的——而是深深打入地下的。中国尊（528米）的桩基础最深达到地下约100米。这些桩就像大树的根，把整栋楼牢牢"钉"在坚实的岩层上。

\textbf{抗震：}北京在历史上发生过多次地震（1679年三河-平谷8级地震、1976年唐山7.8级地震对北京也有影响）。所以北京的高楼都有"隔震支座"——这是一种特殊的橡胶-钢板叠层垫，放在大楼的底部。当地震发生时，隔震支座让大楼"滑动"而不是"硬扛"——可以把地震力降低50-70\%。

\begin{figplaceholder}
插图：隔震支座的剖面——橡胶层+钢板层交错叠放，在地震时产生剪切变形
\end{figplaceholder}

\textbf{风振控制：}在330米的高度，风可能让大楼产生幅度约半米的摇晃——住在顶层的居民会感到"晕船"。中国尊在顶部装了一个重达1000吨的"调谐质量阻尼器"（TMD）——它是一个巨大的摆锤。当大楼被风吹向一边时，摆锤向反方向摆动——抵消了晃动。

\begin{thinkbox}
你坐电梯到一栋高楼的顶层时，有没有感觉大楼在微风中晃动？为什么大楼可以"摇"但不能"倒"？
\end{thinkbox}

\begin{funfact}
上海中心大厦（632米）顶部的"上海慧眼"阻尼器重达1000吨——是目前世界上最大的阻尼器。当台风把大楼吹向左侧时，这个千斤重的"大摆锤"同时向右侧摆回来——游客在顶层根本感觉不到大楼在晃动。
\end{funfact}


\newpage

\section{大兴机场——全球最大单体航站楼是怎么建成的？}

2019年9月，北京大兴国际机场投运——全球最大的单体航站楼，航站楼面积70万平方米（相当于98个足球场）。从开工到投运只用了约4年——这个速度在全世界都是罕见的。

最震撼的是它的屋顶——一个巨大的自由曲面钢网壳，覆盖面积18万平方米，由超过6万根钢管和2万多个焊接球节点组成。

\begin{figplaceholder}
插图：大兴机场的屋顶钢结构——钢网壳的节点连接示意
\end{figplaceholder}

\textbf{几个让人震惊的数字：}
\begin{itemize}
  \item 钢结构总重量约1.4万吨——相当于两个埃菲尔铁塔
  \item 最大跨度约200米——中间没有柱子
  \item 屋顶的曲面由13个"C形柱"支撑——C形柱同时承担支撑、采光、排水的功能
\end{itemize}

\textbf{它是怎么吊装的？}——钢结构在工厂分片预制，运到现场后在地面拼装成"大块"，然后用多台大型履带吊车联合吊装——每块"拼图"重达数百吨。

\begin{thinkbox}
大兴机场的屋顶如果下雨，水怎么排？——因为屋顶是双曲面的，雨水会自动流向特定的排水点。这就是"建筑几何"的妙用——形状本身就解决了排水问题，不需要额外的排水管道。
\end{thinkbox}

\begin{funfact}
大兴机场的停车楼有4321个车位——里面有机器人代泊车系统。你把车开到一个指定的"交接区"，下车，机器人从底下把车"抬起来"，送到停车位。取车时在APP上预约——机器人5分钟内把你的车送到出口。
\end{funfact}

\end{document}
